\documentclass[a4paper,10pt]{article}
\usepackage{subfiles}
\usepackage[margin=1in]{geometry}
\usepackage{adjustbox}
\usepackage{amsfonts}
\usepackage{amsmath}
\usepackage{amssymb}
\usepackage{amsthm}
\usepackage{calc}
\usepackage{enumitem}
\usepackage{eucal}
\usepackage{fancyhdr}
\usepackage{mathrsfs}
\usepackage{mathtools}
\usepackage{nameref}
\usepackage{parskip}
\usepackage{physics}
\usepackage{setspace}
\usepackage{tabularx}
\usepackage[most]{tcolorbox}
\usepackage{tikz}
\usepackage{tikz-cd}
\usepackage{titlesec}
\usepackage{verbatim}
\usepackage{xcolor}
\usepackage{quiver}
\usepackage{imakeidx}
\usepackage{hyperref}


\usetikzlibrary{calc}
\usetikzlibrary{decorations.pathmorphing}
\usetikzlibrary{spath3}

% A TikZ style for curved arrows of a fixed height, due to AndreC.
\tikzset{curve/.style={settings={#1},to path={(\tikztostart)
    .. controls ($(\tikztostart)!\pv{pos}!(\tikztotarget)!\pv{height}!270:(\tikztotarget)$)
    and ($(\tikztostart)!1-\pv{pos}!(\tikztotarget)!\pv{height}!270:(\tikztotarget)$)
    .. (\tikztotarget)\tikztonodes}},
    settings/.code={\tikzset{quiver/.cd,#1}
        \def\pv##1{\pgfkeysvalueof{/tikz/quiver/##1}}},
    quiver/.cd,pos/.initial=0.35,height/.initial=0}

% A TikZ style for shortening paths without the poor behavior of `shorten <' and `shorten >'.
\tikzset{between/.style n args={2}{/tikz/execute at end to={
    \tikzset{spath/split at keep middle={current}{#1}{#2}}
}}}

% TikZ arrowhead/tail styles.
\tikzset{tail reversed/.code={\pgfsetarrowsstart{tikzcd to}}}
\tikzset{2tail/.code={\pgfsetarrowsstart{Implies[reversed]}}}
\tikzset{2tail reversed/.code={\pgfsetarrowsstart{Implies}}}
% TikZ arrow styles.
\tikzset{no body/.style={/tikz/dash pattern=on 0 off 1mm}}

%%%%%%%%%%%%%%%%%%%%%%%%%%%%%%%%%%%%%%%%%%%%%%%%%%%%%%%%%%%%%%%%%%%%%%%%%%%%%%%%

\hypersetup{
    colorlinks,
    citecolor=black,
    filecolor=black,
    linkcolor=blue!60!black,
    urlcolor=blue!60!black
}
\makeindex[title=Index]
\makeindex[name=thm,title=Index of Theorems]


\titleformat{name=\section}
{\normalfont\Large\bfseries}
{}
{0pt}
{}

\tcbuselibrary{theorems}
\NewTcbTheorem[number within=section]{theorem}{Theorem}{
colback=white,
colframe=gray!30,
fonttitle=\bfseries,
coltitle=black,
colbacktitle=gray!30,
toptitle=1mm,
bottomtitle=1mm,
parbox=false}{tcb}
\NewTcbTheorem[use counter from=theorem]{lemma}{Lemma}{
colback=white,
colframe=gray!30,
fonttitle=\bfseries,
coltitle=black,
colbacktitle=gray!30,
toptitle=1mm,
bottomtitle=1mm,
parbox=false}{tcb}
\NewTcbTheorem[use counter from=theorem]{exercise}{Exercise}{
colback=white,
colframe=gray!30,
fonttitle=\bfseries,
coltitle=black,
colbacktitle=gray!30,
toptitle=1mm,
bottomtitle=1mm,
parbox=false}{tcb}
\NewTcbTheorem[use counter from=theorem]{corollary}{Corollary}{
colback=white,
colframe=gray!30,
fonttitle=\bfseries,
coltitle=black,
colbacktitle=gray!30,
toptitle=1mm,
bottomtitle=1mm,
parbox=false}{tcb}
\NewTcbTheorem[use counter from=theorem]{proposition}{Proposition}{
colback=white,
colframe=gray!30,
fonttitle=\bfseries,
coltitle=black,
colbacktitle=gray!30,
toptitle=1mm,
bottomtitle=1mm,
parbox=false}{tcb}

\theoremstyle{definition}
\newtheorem{remark}[tcb@cnt@theorem]{Remark}
\newtheorem{example}[tcb@cnt@theorem]{Example}
\newtheorem{warning}[tcb@cnt@theorem]{Warning}
\newtheorem{notation}[tcb@cnt@theorem]{Notation}
\newtheorem{definition}[tcb@cnt@theorem]{Definition}
\newtheorem{motivation}[tcb@cnt@theorem]{Motivation}
\newtheorem{application}[tcb@cnt@theorem]{Application}
\newtheorem{construction}[tcb@cnt@theorem]{Construction}

\newcommand\aref[2][blue!60!black]{\begingroup\hypersetup{linkcolor=#1}\hyperlink{#2}{#2}\endgroup}
\newcommand\bref[3][blue!60!black]{\begingroup\hypersetup{linkcolor=#1}\hyperlink{#2}{#3}\endgroup}
\newcommand\lref[2][blue!60!black]{lemma \ref{tcb:#2}}
\newcommand\tref[2][blue!60!black]{theorem \ref{tcb:#2}}
\newcommand\eref[2][blue!60!black]{exercise \ref{tcb:#2}}
\newcommand\cref[2][blue!60!black]{corollary \ref{tcb:#2}}
\newcommand\pref[2][blue!60!black]{proposition \ref{tcb:#2}}
\newcommand\Lref[2][blue!60!black]{Lemma \ref{tcb:#2}}
\newcommand\Tref[2][blue!60!black]{Theorem \ref{tcb:#2}}
\newcommand\Eref[2][blue!60!black]{Exercise \ref{tcb:#2}}
\newcommand\Cref[2][blue!60!black]{Corollary \ref{tcb:#2}}
\newcommand\Pref[2][blue!60!black]{Proposition \ref{tcb:#2}}

\makeatletter
\newcommand*{\shortenfences}[3]{%
  % #1: left fence without \left
  % #2: formula inside the fences
  % #3: right fence without \right
  \mathpalette{\@shortenfences{#1}{#3}}{#2}%
}
\newdimen\sf@dimen
\newcommand*{\@shortenfences}[4]{%
  % #1: left fence
  % #2: right fence
  % #3: math style
  % #4: formula
  \sbox0{$#3#4\m@th$}%
  \sbox2{$#3\vcenter{}$}%
  % \dimen0: height above math axis
  \dimen0=\dimexpr\ht0 - \ht2\relax
  \ifdim\dimen0<\z@
    \dimen0=\z@
  \fi
  % \dimen2: depth below math axis
  \dimen2=\dimexpr\ht2 + \dp0\relax
  \ifdim\dimen2<\z@
    \dimen2=\z@
  \fi
  % \sf@dimen: amount for lowering the inner formula
  % to center the inner formula.
  \sf@dimen=\dimexpr(\dimen0 - \dimen2)/2\relax
  % lower the inner formula and raise the outer formula with
  % the fences to keep the math axis.
  \raisebox{\sf@dimen}{%
    $#3\left#1\raisebox{-\sf@dimen}{\box0}\right#2\m@th$%
  }
}
\makeatother

\newcommand\eq{{\operatorname{eq}}}
\newcommand\GL{{\operatorname{GL}}}
\newcommand\Gr{{\operatorname{Gr}}}
\newcommand\id{{\operatorname{id}}}
\newcommand\pr{{\operatorname{pr}}}
\newcommand\pt{{\operatorname{pt}}}
\newcommand\sk{{\operatorname{sk}}}
\newcommand\SL{{\operatorname{SL}}}
\newcommand\cod{{\operatorname{cod}}}
\newcommand\dom{{\operatorname{dom}}}
\newcommand\fin{{\operatorname{fin}}}
\newcommand\Fun{{\operatorname{Fun}}}
\newcommand\Gal{{\operatorname{Gal}}}
\newcommand\Hom{{\operatorname{Hom}}}
\newcommand\lex{{\operatorname{lex}}}
\newcommand\Map{{\operatorname{Map}}}
\newcommand\Mat{{\operatorname{Mat}}}
\newcommand\Mor{{\operatorname{Mor}}}
\newcommand\qis{{\operatorname{qis}}}
\newcommand\red{{\operatorname{red}}}
\newcommand\rev{{\operatorname{rev}}}
\newcommand\cart{{\operatorname{cart}}}
\newcommand\coeq{{\operatorname{coeq}}}
\newcommand\Flag{{\operatorname{Flag}}}
\newcommand\Isom{{\operatorname{Isom}}}
\newcommand\Bilin{{\operatorname{Bilin}}}
\newcommand\canon{{\operatorname{canon}}}
\newcommand\revlex{{\operatorname{revlex}}}


\newcommand\Bl{\operatorname{Bl}}
\newcommand\Cl{\operatorname{Cl}}
\newcommand\gr{\operatorname{gr}}
\newcommand\im{\operatorname{im}}
\newcommand\pd{\operatorname{pd}}
\newcommand\adj{\operatorname{adj}}
\newcommand\Ann{\operatorname{Ann}}
\newcommand\Ass{\operatorname{Ass}}
\newcommand\Aut{\operatorname{Aut}}
\newcommand\Der{\operatorname{Der}}
\newcommand\Div{\operatorname{div}}
\newcommand\DIV{\operatorname{Div}}
\newcommand\Eig{\operatorname{Eig}}
\newcommand\End{\operatorname{End}}
\newcommand\Ext{\operatorname{Ext}}
\newcommand\Fix{\operatorname{Fix}}
\newcommand\kod{\operatorname{kod}}
\newcommand\Kom{\operatorname{Kom}}
\newcommand\lcm{\operatorname{lcm}}
\newcommand\Nil{\operatorname{Nil}}
\newcommand\Orb{\operatorname{Orb}}
\newcommand\Pic{\operatorname{Pic}}
\newcommand\rad{\operatorname{rad}}
\newcommand\rem{\operatorname{rem}}
\newcommand\Syl{\operatorname{Syl}}
\newcommand\Sym{\operatorname{Sym}}
\newcommand\syz{\operatorname{syz}}
\newcommand\Tor{\operatorname{Tor}}
\newcommand\Bord{\operatorname{Bord}}
\newcommand\CaCl{\operatorname{CaCl}}
\newcommand\coim{\operatorname{coim}}
\newcommand\cone{\operatorname{cone}}
\newcommand\Curl{\operatorname{curl}}
\newcommand\Disc{\operatorname{Disc}}
\newcommand\Frac{\operatorname{Frac}}
\newcommand\Grad{\operatorname{grad}}
\newcommand\mult{\operatorname{mult}}
\newcommand\Proj{\operatorname{Proj}}
\newcommand\sign{\operatorname{sign}}
\newcommand\Spec{\operatorname{Spec}}
\newcommand\Stab{\operatorname{Stab}}
\newcommand\Supp{\operatorname{Supp}}
\newcommand\Triv{\operatorname{Triv}}
\newcommand\coker{\operatorname{coker}}
\newcommand\codim{\operatorname{codim}}
\newcommand\colim{\operatorname{colim}}
\newcommand\depth{\operatorname{depth}}
\newcommand\graph{\operatorname{graph}}
\newcommand\MSpec{\operatorname{MSpec}}
\newcommand\Slice{\operatorname{Slice}}
\newcommand\OrBord{\operatorname{OrBord}}
\newcommand\height{\operatorname{height}}

\newcommand\NL{\mathbb{NL}}
\newcommand\trdeg{\operatorname{tr.deg}}
\newcommand\sext{{\operatorname{\mathscr{E}\hspace{-1pt}xt}}}
\newcommand\shom{{\operatorname{\mathscr{H}\hspace{-1.5pt}om}}}
\newcommand\tcap{\mathbin
{\mathmakebox[\widthof{$\cap$}][c]{\top}\mathmakebox[0pt][r]{\cap}}}
\newcommand\pp[2][]{\frac{\partial#1}{\partial#2}}
\newcommand\ppn[3][]{\frac{\partial^{#3}#1}{\partial#2^{#3}}}
\newcommand\ppe[3][]{{\mathmakebox[\widthof{$\displaystyle\frac{\partial#1}
{\partial#2}|$}][r]{\left.\frac{\partial#1}{\partial#2}\right|}}_{#3}}
\newcommand\ppne[4][]{{\mathmakebox[\widthof{$\displaystyle\frac{\partial#1^{#3}}
{\partial#2^{#3}}|$}][r]{\left.\frac{\partial#1}{\partial#2}\right|}}_{#4}}
\newcommand\tpp[2][]{\partial#1/\partial#2}
\newcommand\dpp[2][]{\displaystyle\frac{\partial#1}{\partial#2}}

\renewcommand\op{{\operatorname{op}}}

\DeclareFontFamily{U}{mathx}{}
\DeclareFontShape{U}{mathx}{m}{n}{<-> mathx10}{}
\DeclareSymbolFont{mathx}{U}{mathx}{m}{n}
\DeclareMathAccent{\widecheck}{0}{mathx}{"71}

\DeclareSymbolFont{bbold}{U}{bbold}{m}{n}
\DeclareSymbolFontAlphabet{\mathbbold}{bbold}
\DeclareMathSymbol{\bbomega}{\mathord}{bbold}{"7F}
\DeclareMathSymbol{\bbk}{\mathord}{bbold}{"6B}
\newcommand\ncat[1]{\mathbbold{#1}}

\newcommand\di[1]{\index{#1}\textbf{\textcolor{orange}{#1}}}
\newcommand\dni[1]{\textbf{\textcolor{orange}{#1}}}
\newcommand\cbox[1]{\colorbox{gray!20}{#1}}
\newcommand\dbox[1]{\colorbox{blue!10}{#1}}
\newcommand\ebox[1]{\colorbox{orange!10}{#1}}
\newcommand\cbl[1]{\colorbox{red!10}{#1}}
\newcommand\inj{\mathrel{\hookrightarrow}}
\newcommand\surj{\mathrel{\twoheadrightarrow}}
\newcommand\ol[1]{\overline{#1}}
\newcommand\ul[1]{\underline{#1}}
\newcommand\wh[1]{\widehat{#1}}
\newcommand\wt[1]{\widetilde{#1}}
\newcommand\wc[1]{\widecheck{#1}}
\renewcommand\leq{\leqslant}
\renewcommand\geq{\geqslant}
\renewcommand\mod[1]{\ (\mathrm{mod}\ #1)}
\renewcommand\preceq{\preccurlyeq}
\renewcommand\succeq{\succcurlyeq}
\newcommand\cat[1]{{\textsf{#1}}}
\newcommand\lsub[2]{\vphantom{#2}_{#1}{#2}}
\newcommand\lsup[2]{\vphantom{#2}^{#1}{#2}}

\newcommand\circled[1]{\tikz[baseline=(char.base)]{
\node[shape=circle,draw,inner sep=2pt](char){#1};}}

\newcommand\pagebreaktoggle{}
\newcommand\nonmath[1]{}

\newcommand\sproj[1]{[Stacks Project \href{https://stacks.math.columbia.edu/tag/#1}{#1}]}
\newcommand\krd[1]{[Kerodon \href{https://kerodon.net/tag/#1}{#1}]}
\newcommand\mbook[1]{[Moduli Book \href{https://book.themoduli.space/tag/#1}{#1}]}
\newcommand\mse[1]{[MSE \href{https://math.stackexchange.com/questions/#1}{#1}]}
\newcommand\mo[1]{[MO \href{https://mathoverflow.net/questions/#1}{#1}]}
\newcommand\arxiv[1]{[arXiv \href{https://arxiv.org/abs/math/#1}{#1}]}

\raggedbottom
\pagestyle{fancy}
\lhead{\bref{toc}{ToC}\quad\bref{index}{Index}}
\chead{}
\rhead{}
\cfoot{\thepage}

\setlength{\headheight}{15pt}
\addtolength{\topmargin}{-3pt}
\setlength{\parindent}{16pt}
\setlength{\parskip}{6pt}
\makeatletter
\g@addto@macro\normalsize{%
  \setlength\abovedisplayskip{6pt}
  \setlength\belowdisplayskip{6pt}
  \setlength\abovedisplayshortskip{3pt}
  \setlength\belowdisplayshortskip{3pt}
}
\makeatother
\setlist{listparindent=\parindent,parsep=\parsep,itemsep=0pt,topsep=0pt,partopsep=0pt}
\renewcommand\arraystretch{1}
\renewcommand\baselinestretch{1.2}
\allowdisplaybreaks
\emergencystretch 3em
\begin{document}
\def\numberline#1{}
\hypertarget{toc}{}
\tableofcontents
\pagebreak

\section{Composition Series for Primary Ideals}
\begin{notation}
All rings will be unital and commutative.
\end{notation}

\begin{theorem}{}{thm1}
Let $R$ be a ring. 
The following are equivalent:
\begin{enumerate}[label=(\arabic*)]
\item 
$R$ is Noetherian and every prime ideal of $R$ is maximal.
\item
$R$ is Noetherian and $\dim R=0$.
\item
$R$ is of finite length as an $R$-module.
\item 
$R$ is Artinian.
\end{enumerate}
\end{theorem}
\begin{proof}
Eisenbud Theorem 2.14.
\end{proof}

\begin{corollary}{}{cor1}
Let $R$ is a Noetherian ring.
The following are equivalent:
\begin{enumerate}[label=(\arabic*)]
\item 
Every prime ideal of $R$ is maximal.
\item
$\dim R=0$.
\item
$R$ is of finite length as an $R$-module.
\item 
$R$ is Artinian.
\end{enumerate}
\end{corollary}
\begin{proof}
Follows from \tref{thm1} and the fact that Artinian rings are Noetherian.
\end{proof}

\begin{corollary}{}{cor2}
Let $R$ be a Noetherian ring and $\mathfrak{q}$ be a 
$\mathfrak{p}$-primary ideal of $R$.
Then $R/\mathfrak{q}$ is of finite length as an $R/\mathfrak{q}$-module 
if and only if $\mathfrak{p}$ is a maximal ideal of $R$.

Moreover, any $R$-submodule of $R/\mathfrak{q}$ is 
a $R/\mathfrak{q}$-submodule of $R/\mathfrak{q}$.
So $R/\mathfrak{q}$ is of finite length as an $R/\mathfrak{q}$-module 
if and only if $R/\mathfrak{q}$ is of finite length as an $R$-module.
\end{corollary}
\begin{proof}
$R$ is Noetherian implies $R/\mathfrak{q}$ is Noetherian.
By \cref{cor1}, $R/\mathfrak{q}$ 
is of finite length as an $R/\mathfrak{q}$-module if and only if
every prime ideal of $R/\mathfrak{q}$ is maximal,
which happens if and only if $\mathfrak{p}$ is a maximal ideal of $R$
(because prime ideals of $R/\mathfrak{q}$ are in bijection 
with prime ideal of $R$ containing $\mathfrak{q}$).

An $R$-submodule of $R/\mathfrak{q}$ is the quotient of a $R$-submodule of $R$
containing the $R$-module $\mathfrak{q}$ by the $R$-module $\mathfrak{q}$.
So it is of the form $I/\mathfrak{q}$ for some ideal $I\subseteq R$ 
containing $\mathfrak{q}$.
Since $\mathfrak{q}$ annihilate $I/\mathfrak{q}$, it follows that 
$I/\mathfrak{q}$ can be viewed as a $R/\mathfrak{q}$.
Conversely, any $R/\mathfrak{q}$-submodule of $R/\mathfrak{q}$
can be viewed as an $R$-submodule of $R/\mathfrak{q}$ by the 
surjective ring homomorphism $R\to R/\mathfrak{q}$.
\end{proof}

\begin{lemma}{}{lem1}
Let $R$ be a ring.
For any ideal $I\subseteq R$ and element $x\in R$, 
the following is a surjective $R$-module homomorphism:
\begin{equation*}\begin{array}{r@{}c@{}l}
R&{}\longrightarrow{}&(I+(x))/I\\
1&{}\longmapsto{}&x+I\\
a&{}\longmapsto{}&ax+I
\end{array}\end{equation*}
with kennel $(I:x)=\{a\in R\mid ax\in I\}$.
In particular, there is an $R$-module isomorphism 
$R/(I:x)\cong(I+(x))/I$.
\end{lemma}
\begin{proof}
The $R$-module homomorphism $R\to(I+(x))/I$ defined by $1\mapsto x+I$ is 
surjective because $(I+(x))/I$ is generated by $x+I$ as an $R$-module
(the coset of $(I+(x))/I$ represented by $i+ax\in I+(x)$ is the 
equal to the coset represented by $ax$ which is the image of $a\in R$).

The kernel of this $R$-module homomorphism is by definition 
\begin{equation*}
\{a\in R\mid ax+I=0+I\text{ in }(I+(x))/I\}=\{a\in R\mid ax\in I\}=(I:x)\qedhere
\end{equation*}
\end{proof}

\begin{corollary}{}{cor3}
Let $R$ be a ring.
For any ideal $I\subseteq R$ and element $x\in R$,
there is an bijection between the set of ideals $J\subseteq R$ 
such that $I\subseteq J\subseteq I+(x)$ and the set of ideals
of $R$ containing $(I:x)$.
Moreover, the bijection is given by $J\mapsto I+xJ$.
\end{corollary}
\begin{proof}
This follows from the correspondence theorem of submodules applied to the 
isomorphism $R/(I:x)\cong(I+(x))/I$ in \lref{lem1}.
The isomorphism is given by $1\mapsto x$, 
so an ideal $J/(I:x)$ is mapped to the submodule 
\begin{equation*}
\frac{I+xJ}{I}\subseteq\frac{I+(x)}{I}=\frac{I+xR}{I}\qedhere
\end{equation*}
\end{proof}

\begin{application}
Let $k$ be a field and $R=k[z_0,\dots,z_{c-1}]$ be a polynomial ring over $k$.
Let $\mathfrak{m}=(x_0,\dots,x_{n-1})$ be a maximal ideal of $R$
and $I=(f_0,\dots,f_{r-1})$ be an $\mathfrak{m}$-primary ideal of $R$.
The weird indexing is to match the indexing of the variables in Macaulay2.

The following is a algorithm to compute a maximal ascending chain of ideals 
$I=I_0\subsetneq I_1\subsetneq\cdots\subsetneq I_s=\mathfrak{m}$
containing $I$ up to the maximal ideal $\mathfrak{m}$.
Taking the quotients of $R$ by this chain gives a composition series of $R/I$.
\begin{itemize}
\item
Define $M_i:=I+(x_0,\dots,x_{i-1})$ for each $i\in\{0,\dots,n\}$.
\item
Define $Q_i:=M_{i+1}/M_i$ for each $i\in\{0,\dots,n-1\}$.
\item 
Define $J_i:=(M_i:x_i)$ for each $i\in\{0,\dots,n-1\}$.
\end{itemize}
We begin the output chain with $I=M_0$ and for each $i\in\{0,\dots,n-1\}$, 
we attach a maximal ascending chain of ideal of $R$ in between $M_i$ and 
$M_{i+1}$ (excluding $M_i$ and including $M_{i+1}$) to the output chain.
There are a three cases:
\begin{enumerate}[label=Case \arabic*.,align=left]
\item
If $Q_i=0$, then $M_{i+1}=M_i$.
So we do not attach any ideals to the output chain.
\item
If $Q_i$ is simple, then we attach $M_{i+1}$ to the output chain.
\item 
If $Q_i$ is not simple, we invoke the isomorphism 
\begin{equation*}
R/J_i\cong Q_i=M_{i+1}/M_i
\end{equation*}
To get a maximal chain of ideals between $M_i$ and $M_{i+1}$,
we compute a maximal chain of ideals between $J_i$ and $R$, 
and for each ideal $K$ inside this chain, we attach 
$M_i+x_iK$ to the output chain (because this is the bijection in \cref{cor3}).

We will not run into an infinite loop because $R/J_i$ necessarity has 
smaller length than $R/I$ as an $R$-module when $Q_i\not=0$.
\end{enumerate}
\end{application}

\section{A Criterion for Ideal of Finite Colength}
\begin{theorem}{Chinese Remainder Theorem}{thm-crt}
Let $R$ be a ring and $I,J\subseteq R$ be ideals.
If $I+J=R$, thenhe injective ring homomorphism 
\begin{equation*}\begin{array}{r@{}c@{}l}
\varphi:R/(I\cap J)&{}\longrightarrow{}&R/I\times R/J\\
a+(I\cap J)&{}\longmapsto{}&(a+I,a+J)
\end{array}\end{equation*}
is also surjective.
In particular, $\varphi$ is an isomorphism
\end{theorem}
\begin{proof}
$I+J=R$ implies there exists $i\in I$ and $j\in J$ such that $1=i+j$.
Then 
\begin{equation*}
\varphi(i+(I\cap J))=(0+I,1+J)\qquad\text{and}\qquad
\varphi(j+(I\cap J))=(1+I,0+J)
\end{equation*}
Thus for any $(a+I,b+J)\in R/I\times R/J$, we have
\begin{equation*}
\varphi(bi+aj+(I\cap J))=(a+I,b+J)\qedhere
\end{equation*}
\end{proof}

\begin{lemma}{}{lem-crt}
Let $R$ be a ring and $I_1,\dots,I_n,J\subseteq R$ be ideals.
Then 
\begin{equation*}
(I_1\cap\cdots\cap I_n)+J\subseteq(I_1+J)\cap\cdots\cap(I_n+J)
\end{equation*}
Moreover, if $I_1+J=\cdots=I_n+J=R$, the equality holds.
\end{lemma}
\begin{proof}
$(I_1\cap\cdots\cap I_n)+J\subseteq(I_1+J)\cap\cdots\cap(I_n+J)$ is clear.
Suppose $I_1+J=\cdots=I_n+J=R$.
Then for each $i\in\{1,\dots,n\}$, there exists $a_i\in I_i$ and $b_i\in J$
such that $1=a_i+b_i$.
If follows that
\begin{equation*}
1=(a_1+b_1)\cdots(a_n+b_n)=(a_1\cdots a_n)+
\underbrace{\cdots}_{\mathmakebox[0pt][c]
{\text{all remaining term has a factor of some $b_i$, so it is in $J$}}}
\in(I_1\cap\cdots\cap I_n)+J.
\end{equation*}
So $(I_1\cap\cdots\cap I_n)+J=R=(I_1+J)\cap\cdots\cap(I_n+J)$.
\end{proof}

\begin{corollary}{Chinese Remainder Theorem}{cor-crt}
Let $R$ be a ring and $I_1,\dots,I_n\subseteq R$ be ideals
such that $I_i+I_j=R$ for all distinct $i,j\in\{1,\dots,n\}$.
Then 
\begin{equation*}
\frac{R}{I_1\cap\cdots\cap I_n}\cong\frac{R}{I_1}\times\cdots\times\frac{R}{I_n}
\end{equation*}
\end{corollary}
\begin{proof}
We have the chain of equalities by \tref{thm-crt}:
\begin{equation*}
\frac{R}{I_1\cap\cdots\cap I_n}\cong\frac{R}{I_1\cap\cdots\cap I_{n-1}}\times\frac{R}{I_n}
\cong\cdots\cong\frac{R}{I_1\cap I_2}\times\frac{R}{I_3}\times\cdots\times\frac{R}{I_n}
\cong\frac{R}{I_1}\times\cdots\times\frac{R}{I_n}
\end{equation*}
For the $i$-th isomorphism, the assumption we need to verify in order to 
invoke \tref{thm-crt} is that $(I_1\cap\cdots\cap I_{i-1})+I_i=R$.
This follows from \lref{lem-crt} and the assumption that 
$I_j+I_i=R$ for all $j\in\{1,\dots,i-1\}$.
\end{proof}

\begin{lemma}{}{lem4}
Let $R$ be a ring and $\mathfrak{q}_1$ be a $\mathfrak{p}_1$-primary ideal
and $\mathfrak{q}_2$ be a $\mathfrak{p}_2$-primary ideal of $R$.
If $\mathfrak{p}_1+\mathfrak{p}_2=R$, then $\mathfrak{q}_1+\mathfrak{q}_2=R$. 
\end{lemma}
\begin{proof}
Suppose $\mathfrak{p}_1+\mathfrak{p}_2=R$, then there exists 
$a\in\mathfrak{p}_1$ and $b\in\mathfrak{p}_2$ such that $a+b=1$.
Since $\mathfrak{p}_1=\sqrt{\mathfrak{q}_1}$ and 
$\mathfrak{p}_2=\sqrt{\mathfrak{q}_2}$,
it follows that there exists $m,n\in\mathbb{Z}_{>0}$ such that
$a^m\in\mathfrak{q}_1$ and $b^n\in\mathfrak{q}_2$.
Then 
\begin{equation*}
1=(a+b)^{m+n}\in\mathfrak{q}_1+\mathfrak{q}_2\qedhere
\end{equation*}
\end{proof}

\begin{theorem}{}{thm2}
Let $R$ be a noetherian ring and $I$ be a proper ideal
so that $I$ admit a primary decomposition 
$I=\mathfrak{q}_1\cap\cdots\cap\mathfrak{q}_n$.
Let $\MSpec(R)=\{\text{maximal ideals of }R\}$ and 
$\Ass(I)=\{\mathfrak{p}_1\cap\cdots\cap\mathfrak{p}_n\}$
where $\mathfrak{p}_i=\sqrt{\mathfrak{q}_i}$.

Then $R/I$ is of finite length as an $R$-module if and only if 
the $\Ass(I)\subseteq\MSpec(R)$.
\end{theorem}
\begin{proof}
Suppose $\Ass(I)\subseteq\MSpec(R)$.
Then each $\mathfrak{p}_i$ is a maximal ideal and 
$\mathfrak{p}_i+\mathfrak{p}_j=R$ for distinct $i,j\in\{1,\dots,n\}$
since maximal ideals are pairwise comaximal.

Then by \lref{lem4}, we have $\mathfrak{q}_i+\mathfrak{q}_j=R$ 
for distinct $i,j\in\{1,\dots,n\}$.
By \cref{cor-crt}, we have
\begin{equation*}
\frac{R}{I}=\frac{R}{\mathfrak{q}_1\cap\cdots\cap\mathfrak{q}_n}
\cong\frac{R}{\mathfrak{q}_1}\times\cdots\times\frac{R}{\mathfrak{q}_n}
\end{equation*}
By \cref{cor2}, each $R/\mathfrak{q}_i$ is of finite length as $R$-modules 
because $\mathfrak{p}_i=\sqrt{\mathfrak{q}_i}$ is a maximal ideal.
Product of finite length modules is of finite length,
so $R/I$ is of finite length as an $R$-module.

Conversely, suppose $\Ass(I)\not\subseteq\MSpec(R)$.
Then there exists $i\in\{1,\dots,n\}$ such that $\mathfrak{p}_i$ 
is not a maximal ideal.
Now $I=\mathfrak{q}_1\cap\cdots\cap\mathfrak{q}_n\subseteq
\mathfrak{q}_i\subseteq\mathfrak{p}_n$ which implies 
there exists a surjective ring homomorphism $R/I\to R/\mathfrak{p_i}$.
Since $\mathfrak{p}_i$ is a prime ideal but not a maximal ideal,
It follows that $\dim R/\mathfrak{p}_i\geq1$.
Then $\dim R/I\geq1$ as well which implies 
$R/I$ is not of finite length as an $R/I$-module by \cref{cor1}.
By \cref{cor2}, $R/I$ is not of finite length as an $R$-module.
\end{proof}

\pagebreak

\printindex
\hypertarget{index}{}
\phantomsection
\addcontentsline{toc}{section}{Index}
\end{document}